\documentclass[11pt]{article}
\usepackage[margin=0.9in]{geometry}
\usepackage{amsmath,amssymb,amsthm,booktabs}
\usepackage[T1]{fontenc}
\usepackage{lmodern,microtype}
\usepackage[colorlinks=true,linkcolor=blue,urlcolor=blue]{hyperref}
\newtheorem{theorem}{Theorem}[section]
\newtheorem{lemma}[theorem]{Lemma}
\newtheorem{proposition}[theorem]{Proposition}
\newtheorem{corollary}[theorem]{Corollary}
\theoremstyle{remark}
\newtheorem{remark}[theorem]{Remark}
\newcommand{\R}{\mathbb R}
\newcommand{\C}{\mathbb C}
\newcommand{\norm}[1]{\left\|#1\right\|}
\newcommand{\diag}{\operatorname{diag}}
\newcommand{\ran}{\operatorname{ran}}
\newcommand{\HH}{\mathcal H}
\title{Toward a four-dimensional Yang--Mills construction\\
\large Notebook 08: a local flat minimizing extension}
\author{Research working document}
\date{16 October 2026}
\setlength{\emergencystretch}{5em}
\begin{document}
\maketitle

\begin{abstract}
We construct an explicit minimizing right inverse for the homogeneous
flat Maxwell block average. A scalar gauge choice with a squared
Laplacian denominator removes the remaining complex massless poles
from the minimizing field itself. The resulting potential and every
first fine difference have exponential locality in coarse-cell
distance, uniformly in refinement. The potential has physical
discrete $C^{1,\alpha}$ bounds and its curvature has $C^\alpha$
bounds for every fixed $\alpha<1$. The scalar gauge matches the
projection in the deposited homogeneous flat source formulas.
The componentwise fine Laplacian is also uniformly bounded and local;
this uses the physical averaging adjoint rather than an absolute
second-derivative Fourier sum. The
corrected bond average and its normalization are checked separately.
This supplies local flat response legs, but does not prove a
nonflat or independent-current analytic response theorem.
\end{abstract}

\noindent\textbf{Status.}
Research proof checked by separate internal derivations and audits.
No external certification or novelty claim is made.
The explicit quadratic construction is proved
below. The final source identification uses deposited transcriptions
of the primary operator definitions and their stated minimizing
property; the unavailable original CMP99b proof has not been
independently reverified here.

\section{The minimizing problem and the local norms}

Let $M$ be odd, $\eta=M^{-1}$, and use centered $M$-blocks.
The normalized straight-string average and fine derivative are
\[
 (Q_M A)_\mu(n)=M^{-d-1}\sum_{s\in I_M^d}
          \sum_{j=0}^{M-1}A_\mu(n+(s+je_\mu)/M),
 \qquad I_M=\{-(M-1)/2,\ldots,(M-1)/2\},
\]
\[
 (D_{\eta,\nu}A)_\mu(x)
   =M[A_\mu(x+e_\nu/M)-A_\mu(x)],\qquad
 (d_\eta A)_{\mu\nu}=D_{\eta,\mu}A_\nu-D_{\eta,\nu}A_\mu.
\]
Physical fine-volume norms have measure $M^{-d}$; unit coarse
norms have counting measure. Define
\[
 q_M(B)=\inf_{Q_M A=B}\norm{d_\eta A}_\eta^2.
\]
Notebook 07 derived its bounded holomorphic Fourier symbol and
uniform exponential locality. We use its notation below and recall
the relevant cancellation explicitly.

For fine values in a fixed physical cell, the local Hilbert norm is
\[
 \norm v_{\mathrm{cell},2}^2
       =M^{-d}\sum_{s\in I_M^d}|v(s/M)|^2.
\]
The physical discrete $C^\alpha$ norm is the supremum norm plus
\[
 \sup_{s\ne t}
 \frac{|v(s/M)-v(t/M)|}{|(s-t)/M|^\alpha}.
\]
The discrete $C^{1,\alpha}$ norm controls $v$, all full fine
differences $D_{\eta,\nu}v$, and the $C^\alpha$ seminorms of
these differences. Use a fixed physical collar whenever a fine
difference or comparison crosses the cell boundary. Constants may
depend on the collar size but never on $M$.

\begin{theorem}\label{thm:extension}
There is a linear extension $\HH_M$ on the infinite lattice and
every compatible periodic lattice such that
\[
 Q_M\HH_M=I,\qquad
 \norm{d_\eta\HH_M B}_\eta^2=q_M(B).
\]
It satisfies the scalar gauge condition defined in Section~\ref{sec:gauge}.
Its reconstructed fine-cell symbol is holomorphic on a coarse
momentum strip independent of $M$. For every fixed $0<\alpha<1$
there are constants $C_\alpha,c_\alpha>0$, independent of $M$,
such that its coarse-cell kernels obey
\begin{equation}\label{eq:locality}
 \begin{split}
 &\norm{\HH_M(n)}_{\C^d\to C^{1,\alpha}_{\mathrm{cell}}}
 +\norm{(d_\eta\HH_M)(n)}_{\C^d\to C^\alpha_{\mathrm{cell}}}\\
 &\qquad+\norm{(\mathcal L_\eta\HH_M)(n)}_{\C^d\to L^\infty_{\mathrm{cell}}}
 \le C_\alpha e^{-c_\alpha|n|_1}.
 \end{split}
\end{equation}
Here $\mathcal L_\eta=\sum_\nu D_{\eta,\nu}^*D_{\eta,\nu}$
acts componentwise on the potential; it is the scalar lattice
Laplacian on each bond component, not the Maxwell curl precision.
The same statement holds in physical cell Hilbert norms.
Periodic kernels are exact periodizations and have uniform bounds
in periodic coarse-cell distance.
\end{theorem}

Colors decouple at flatness; the scalar estimates act on each
orthonormal color component. This theorem is about the stated
full-block flat problem. Its physical source match is specified
only in Section~\ref{sec:source}.

\section{The alias quotient and its canceled symbol}

For $r\in I_M^d$, put $p_r=z+2\pi r$ and define
\[
 a_r=a_M(p_r),\quad
 a_M(p)=\prod_\nu\frac1M\sum_{s\in I_M}e^{ip_\nu s/M},
 \quad
 \omega_{r,\mu}=a_M(p_r)\frac1M\sum_{j=0}^{M-1}e^{ip_{r,\mu}j/M}.
\]
A superscript minus on $a_r$ or $\omega_r$ means evaluation at
$-p_r$, not variable-dependent complex conjugation. Set
\[
 u_{r,\mu}^{\pm}=M(e^{\pm ip_{r,\mu}/M}-1),\quad
 P_r=\sum_\mu u_{r,\mu}^+u_{r,\mu}^-
     =4M^2\sum_\mu\sin^2(p_{r,\mu}/(2M)),
\]
\[
 d_\mu^\pm=e^{\pm iz_\mu}-1,\qquad b_\mu=d_\mu^+d_\mu^-.
\]
Exact string telescoping gives
\begin{equation}\label{eq:cochain}
 \omega_{r,\mu}u_{r,\mu}^+=a_r d_\mu^+.
\end{equation}
The coarse average in a Fourier fiber is
$B_\mu=\sum_r\omega_{r,\mu}A_{r,\mu}$. With physical-volume
Parseval, fine fiber energies sum aliases without an extra $M^d$.

For real nonzero $z$, the fine Maxwell energy in an alias is
\[
 P_r|A_r|^2-|(u_r^+)^*A_r|^2
       =\min_\lambda P_r|A_r+u_r^+\lambda|^2.
\]
Fine gradients average to coarse gradients by \eqref{eq:cochain},
and their scalar average is onto because $a_0\ne0$ on the central
real zone. Thus the minimized form is the gauge quotient of
\[
 W=\diag\left(\sum_r\frac{\omega_{r,\mu}\omega_{r,\mu}^-}{P_r}\right),
\]
namely
\[
 q_M=W^{-1}-W^{-1}d^+(d^-W^{-1}d^+)^{-1}d^-W^{-1}.
\]
At $z=0$, the constant fine field realizes any constant coarse field
with zero energy, so $q_M(0)=0$.

Write
\[
 A=a_0a_0^-,\quad g_\mu=\omega_{0,\mu}\omega_{0,\mu}^-,\quad
 R_\mu=\sum_{r\ne0}\frac{\omega_{r,\mu}\omega_{r,\mu}^-}{P_r},
 \quad D_\mu=g_\mu+P_0R_\mu,\quad h_\mu=D_\mu^{-1},
\]
\[
 h=\diag(h_\mu),\qquad
 E=A^{-1}-\sum_\mu\frac{b_\mu R_\mu}{g_\mu D_\mu}.
\]
The central instance of \eqref{eq:cochain} gives
$\sum_\mu b_\mu/g_\mu=P_0/A$, hence
\begin{equation}\label{eq:q}
 d^-hd^+=P_0E,\qquad
 q_M=P_0h-\frac{hd^+d^-h}{E}.
\end{equation}
This expression has no inverse power of the central massless
denominator $P_0$.

\section{The additional scalar denominator has a uniform domain}

Define
\begin{equation}\label{eq:scalar-two}
 R_{s,2}=\sum_{r\ne0}\frac{a_ra_r^-}{P_r^2},\qquad
 D_{s,2}=A+P_0^2R_{s,2}.
\end{equation}
The scalar weights are
\begin{equation}\label{eq:t}
 t_{2,0}=\frac{a_0^-}{D_{s,2}},\qquad
 t_{2,r}=\frac{P_0^2a_r^-}{P_r^2D_{s,2}}\quad(r\ne0).
\end{equation}
They satisfy $\sum_r a_rt_{2,r}=1$. For real nonzero momentum,
they minimize $\sum_rP_r^2|\lambda_r|^2$ at a prescribed scalar
block average. Indeed the scalar covariance is
$\sum_r a_ra_r^-/P_r^2=D_{s,2}/P_0^2$ and the ordinary weighted
least-squares formula gives \eqref{eq:t}.

\begin{lemma}\label{lem:strip}
There is a complex strip, with a positive real margin around the
central Brillouin cube and width independent of $M$, on which
$A,g_\mu,D_\mu,E,D_{s,2},a_0$ are uniformly separated from zero.
All these functions, their required inverses, $P_0,R_\mu,R_{s,2}$
and $q_M$ are uniformly bounded and holomorphic there.
\end{lemma}

\begin{proof}
Here are the uniform estimates behind the strip argument, so its
use introduces no additional source assumption. Fix $\pi<K<2\pi$
and a small imaginary width $\delta$. For high aliases,
\[
 P_r(x)\ge c_K|r|^2,\qquad
 \Re P_r(x+iy)\ge P_r(x)-e^\delta\sum_\nu y_\nu^2,
 \quad |x_\nu|\le K.
\]
The first bound follows by taking distance from $2\pi M\mathbb Z$;
for a nonzero retained coordinate that distance is at least
$(2\pi-K)|r_\nu|$. Choose $\delta$ so the second term is smaller
than half the high-alias gap. Discrete Parseval gives
\[
 \sum_r|a_r(x+iy)|^2\le e^{\sum|y_\nu|},\quad
 \sum_r|\omega_{r,\mu}(x+iy)|^2
       \le e^{\sum|y_\nu|+2|y_\mu|}.
\]
Cauchy--Schwarz bounds the corresponding bilinear products with
negative arguments. Thus both high remainders in
\eqref{eq:scalar-two} and the vector remainders are uniformly
holomorphic and bounded, despite the growing number of aliases.

On the real central cube, the centered factor
$\sin(t/2)/(M\sin(t/(2M)))$ is between $2/\pi$ and one.
Consequently $A\ge(2/\pi)^{2d}$,
$g_\mu,D_\mu\ge(2/\pi)^{2d+2}$,
$D_\mu\le1+d\pi^2/4$, and
\[
 E\ge\frac{4}{\pi^2(1+d\pi^2/4)}\quad(z\ne0),
 \qquad E(0)=1.
\]
Also $0\le R_{s,2}\le1/16$, so
$D_{s,2}\ge A$. The enlarged real cube has the same estimates
with weaker fixed positive constants: replace $2/\pi$ by
$\sin(K/2)/(K/2)$ and the high gap by
$4(2\pi-K)^2/\pi^2$.

Cauchy estimates on nested open tubes, retaining a real margin,
give uniform derivatives. Integrating in imaginary directions first
keeps $A,g_\mu,D_\mu,D_{s,2}$ nonzero. The formula for $E$ is
then uniformly bounded and holomorphic; a second shrinking and its
real gap keep it nonzero as well. Since $A=a_0a_0^-$ and each
factor is bounded, both factors are bounded away from zero.
All shrinkings are independent of $M$. No Cauchy disk leaves a
previously controlled domain. For $M=1$ the high sums are empty
and the conclusion holds directly.
\end{proof}

\section{The field and its exact central cancellation}

For real nonzero momentum put
\[
 S=d^-hB,\qquad \theta=\frac{S}{P_0E},\qquad
 F_r=P_r^{-1}\diag(\omega_r^-)q_MB.
\]
Define the proposed extension initially by
\begin{equation}\label{eq:uncanceled-H}
 \HH_rB=F_r+u_r^+t_{2,r}\theta.
\end{equation}
This formula explains the gauge choice. Its central term is defined
on the complex null cone by the following cancellation:
\begin{equation}\label{eq:H-central}
 (\HH_0B)_\mu=\omega_{0,\mu}^-h_\mu B_\mu+
 \frac{u_{0,\mu}^+}{a_0}
       \left(R_\mu h_\mu-\frac{P_0R_{s,2}}{D_{s,2}}\right)
       \frac{S}{E}.
\end{equation}
For a high alias,
\begin{equation}\label{eq:H-high}
 \HH_rB=P_r^{-1}\diag(\omega_r^-)q_MB
       +\frac{u_r^+P_0a_r^-S}{P_r^2D_{s,2}E},\qquad r\ne0.
\end{equation}

\begin{proof}[Proof of the cancellation]
The cochain identity and $g_\mu h_\mu=1-P_0R_\mu h_\mu$ give
\[
 \omega_{0,\mu}^-h_\mu d_\mu^+
       =\frac{u_{0,\mu}^+}{a_0}(1-P_0R_\mu h_\mu).
\]
Also $A/D_{s,2}=1-P_0^2R_{s,2}/D_{s,2}$. Substitute these
identities into the central part of \eqref{eq:uncanceled-H}.
The two terms proportional to $P_0^{-1}$ cancel, leaving exactly
\eqref{eq:H-central}. For high aliases, cancel one $P_0$ between
\eqref{eq:t} and $\theta$, obtaining \eqref{eq:H-high}.
\end{proof}

Lemma~\ref{lem:strip} now makes both expressions holomorphic.
At zero, $\HH_0(0)=I$ and every high alias vanishes, so constants
are lifted to constants. No inverse of a zero harmonic eigenvalue
has been introduced.

\section{Constraint, minimal energy and scalar gauge}\label{sec:gauge}

The Feynman minimizing field satisfies
\[
 \sum_r\diag(\omega_r)F_r=B-d^+\theta.
\]
The added gradient averages to $d^+\theta$, by
\eqref{eq:cochain} and $\sum a_rt_{2,r}=1$. Thus $Q_M\HH_M=I$.
Further, $d^-q_M=0$, so
\[
 (u_r^+)^*F_r=(a_r^-/P_r)d^-q_MB=0
\]
on real nonzero fibers. The field $F$ is fine transverse; its
Maxwell and Feynman energies agree and equal $B^*q_MB$.
The added gradient changes neither curl nor energy. This proves
the minimizing identity in Theorem~\ref{thm:extension}, including
the separately specified constant zero mode.

Let $Q_0$ be the scalar block average and
$\Delta=D_f^*D_f$ the positive fine scalar Laplacian. With the
physical volume measures, $Q_0Q_0^*=I$ and $Q_0^*$ extends a
coarse scalar as a blockwise constant fine scalar. For a positive
coarse penalty $a'$, define
\begin{equation}\label{eq:source-R}
 G'=(\Delta+Q_0^*a'Q_0)^{-1},\quad
 P=G'Q_0^*(Q_0G'^2Q_0^*)^{-1}Q_0G',\quad R=I-P.
\end{equation}
On a real finite regulator, $P$ is the orthogonal projection onto
$\ran(G'Q_0^*)$. In particular,
\begin{equation}\label{eq:R-range}
 \ran R=\ker(Q_0G')
         =(\Delta+Q_0^*a'Q_0)\ker Q_0
         =\Delta\ker Q_0.
\end{equation}
On a nonzero Bloch fiber, equivalently
$\ker R=\ran(\Delta^{-1}Q_0^*)$. This follows by taking the
orthogonal complement in \eqref{eq:R-range}; the Laplacian is
invertible in such a fiber. At zero one uses \eqref{eq:source-R},
whose penalty treats scalar constants, rather than an undefined
$\Delta^{-1}$.

The divergence of the extension has the exact form
\begin{equation}\label{eq:divergence}
 (D_f^*\HH_MB)_r
    =P_r^{-1}a_r^-\frac{P_0S}{D_{s,2}E}\qquad(z\ne0).
\end{equation}
The central expression is its canceled value
$a_0^-S/(D_{s,2}E)$. Thus the divergence lies in
$\ran(\Delta^{-1}Q_0^*)$, proving
\begin{equation}\label{eq:gauge}
 R D_f^*\HH_M=0.
\end{equation}
At zero the divergence is zero. The identities extend through the
holomorphic formulas where appropriate; positivity and minimizing
claims are made for real fields.

The gauge is unique among minimizing fields. In a nonzero fiber
the difference of two energy minimizers is a fine gradient
$u_r^+\lambda_r$. Its zero average gives $\sum_ra_r\lambda_r=0$.
Therefore $\Delta\lambda\in\ran R$, while the gauge condition
gives $R\Delta\lambda=0$. Hence the difference vanishes. At zero,
the coarse average fixes the fine constant harmonic one-form and
all other zero-energy differences are high-alias gradients; the
same argument removes them. Equivalently, \eqref{eq:t} minimizes
the divergence norm among the Maxwell energy minimizers. This is
different from minimizing the potential norm itself.

\section{Uniform fine derivative estimates}

The sine bounds used in Lemma~\ref{lem:strip} give, for high aliases
on a fixed smaller tube,
\[
 w(r)=\prod_\nu(1+|r_\nu|)^{-1},\qquad
 |a_r^-|\le Cw(r),\qquad
 |\omega_{r,\mu}^-|\le\frac{Cw(r)}{1+|r_\mu|},
\]
\[
 |u_r^+|\le C(1+|r|),\qquad |P_r|\ge c(1+|r|)^2.
\]
These estimates are valid up to the edge of the fine Brillouin
zone; they do not expand in $r/M$. Equation~\eqref{eq:H-high}
therefore gives
\begin{equation}\label{eq:alias-bound}
 \norm{\HH_r}\le\frac{Cw(r)}{(1+|r|)^2},\qquad
 \norm{u_{r,\nu}^+\HH_r}
       +\norm{u_r^+\wedge\HH_r}
       \le\frac{Cw(r)}{1+|r|}.
\end{equation}
For curl, the gradient term vanishes exactly and one may also use
\[
 u_r^+\wedge\HH_r
       =P_r^{-1}u_r^+\wedge[\diag(\omega_r^-)q_M].
\]
The central aliases and their first fine differences are bounded
directly by \eqref{eq:H-central}.

On a dyadic shell $2^j\le|r|_\infty<2^{j+1}$, choose a largest
coordinate. Its harmonic sum over that shell is bounded, while
each other coordinate contributes at most $C(j+1)$. Thus
\[
 \sum_{\text{shell }j}w(r)\le C_d(j+1)^{d-1}.
\]
Combining with \eqref{eq:alias-bound}, for every $\alpha<1$,
\begin{equation}\label{eq:weighted-alias}
 \sum_r(1+|r|)^\alpha
 \left[\norm{\HH_r}+
       \sum_\nu\norm{u_{r,\nu}^+\HH_r}
       +\norm{u_r^+\wedge\HH_r}\right]\le C_\alpha.
\end{equation}
All constants are uniform in $M$. This absolute sum also proves
uniform alias Hilbert operator bounds. It controls every full first
fine derivative, unlike the simpler least-potential-norm gauge.
This absolute-sum argument does not give the endpoint $\alpha=1$
or pointwise second-derivative bounds. A separate physical-space
identity does give the componentwise Laplacian bound below.

\subsection{The componentwise Laplacian and the physical adjoint}

Multiplication of the uncanceled formula by $P_r$ gives
\begin{equation}\label{eq:laplacian-split}
 P_r\HH_rB=\diag(\omega_r^-)q_MB+\mathcal Z_rB,
 \qquad
 \mathcal Z_rB=\frac{u_r^+a_r^-P_0S}{P_rD_{s,2}E}.
\end{equation}
The central expression is defined by cancellation as
$\mathcal Z_0B=u_0^+a_0^-S/(D_{s,2}E)$.
For high aliases the established bounds give
$\norm{\mathcal Z_r}\le Cw(r)/(1+|r|)$.
The same dyadic summation therefore proves uniform physical
$C^\alpha$ bounds for the reconstructed remainder, on the fixed
complex strip, for every $\alpha<1$.

The first term of \eqref{eq:laplacian-split} is precisely
$Q_M^*q_MB$ with the physical fine-volume adjoint. It should not
be estimated by the absolute sum of its Fourier aliases.
For $x=n+s/M$, $s\in I_M^d$, set
$\beta_s=(s_\mu+(M-1)/2+1)/M$. Direct counting of the strings
through that fine bond gives
\begin{equation}\label{eq:adjoint-row}
 (Q_M^*B)_\mu(n+s/M)
 =\beta_s B_\mu(n)+(1-\beta_s)B_\mu(n-e_\mu),
 \qquad 0\le\beta_s\le1.
\end{equation}
The physical adjoint multiplies the discrete averaging coefficient
by $M^d$. The two string multiplicities sum to $M$, so the two
remaining weights sum to one. Transverse coarse labels are fixed
by the cell containing the fine basepoint. This proves a uniform
fine supremum bound and fixed coarse range for $Q_M^*$.
Its reconstructed Bloch multiplier is
$\beta_s+(1-\beta_s)e^{-iz_\mu}$, uniformly bounded on the strip.
Consequently $Q_M^*q_M$ has uniform supremum-valued holomorphic
cell symbols and exponential coarse-cell locality.

Equation~\eqref{eq:laplacian-split} now proves the claimed
$L^\infty$ bound for $\mathcal L_\eta\HH_M$ and its exponential
locality by the same contour argument used below. The adjoint term
may jump between neighboring transverse blocks, so the proof
does not claim a uniform H\"older bound for the Laplacian or
separate pointwise bounds for all second fine derivatives.

\section{Fine-cell reconstruction and exponential locality}

At a fine position $x=n+s/M$, reconstruct the symbol by
\begin{equation}\label{eq:cell-symbol}
 \HH_{\mathrm{cell}}(z;s)
       =\sum_r e^{i(z+2\pi r)\cdot s/M}\HH_r(z).
\end{equation}
For real $z$, the matrix of exponential factors is unitary from
alias counting-measure $\ell^2$ to the cell norm with measure
$M^{-d}$. For complex $z=x+iy$, multiplication by
$e^{-y\cdot s/M}$ changes this norm and its inverse by at most
$e^{\sum|y_\nu|/2}$. Thus no growing output-dimension factor
appears in the physical Hilbert bound.

For two physical fine positions $\xi,\xi'$ in a fixed cell or
collar, the fixed strip gives
\[
 |e^{ip_r\cdot\xi}-e^{ip_r\cdot\xi'}|
       \le C_\alpha(1+|r|)^\alpha|\xi-\xi'|^\alpha.
\]
This follows from the derivative and bounded-exponential estimates
and $\min(1,t)\le t^\alpha$. Apply
\eqref{eq:weighted-alias} to \eqref{eq:cell-symbol}, and to the
symbols with factors $u_{r,\nu}^+$ or $u_r^+\wedge$.
This proves bounded holomorphic cell symbols in physical discrete
$C^{1,\alpha}$ for the potential and $C^\alpha$ for the curl.
Collars cover differences across cell boundaries. A wrapped fine
point carries the appropriate coarse Bloch factor, exactly as in
the finite-difference symbol $u_r^+$.

The individual alias vector is not a fixed-index periodic symbol.
On real momenta, a shift by $2\pi$ permutes aliases modulo $M$;
the full scalar covariance, vector covariance, quotient and coarse
derivative are periodic. Formula~\eqref{eq:uncanceled-H} therefore
has that permutation law. In \eqref{eq:cell-symbol}, a wrapped
alias changes its phase by $e^{2\pi i s_\nu}=1$.
The reconstructed symbol is genuinely coarse-periodic. The canceled
holomorphic formulas on enlarged momentum charts agree on real
overlaps and hence glue by the identity theorem. This argument
also applies to the fine derivative and curl symbols.

Finally shift each coarse Fourier contour within the common strip.
Uniform operator bounds in the displayed cell norms give
\eqref{eq:locality}, including on fixed collars. The explicit
kernel acts by
\[
 (\HH_MB)(n+s/M)=\sum_m\HH_M(n-m;s)B(m).
\]
Exponential row sums give uniform coarse-$\ell^\infty$ to fine
supremum and local H\"older bounds. The corresponding column sums
use fine physical measure $M^{-d}$, which cancels the $M^d$ fine
positions per cell; weighted Schur gives the coarse-$\ell^2$ to
fine-volume-$\ell^2$ bounds. Finite tori sample the same periodic
symbols, so their kernels are exact periodizations. Absolute
convergence gives uniform periodic-distance estimates, including
small periods with oriented-cell multiplicities.

\section{The homogeneous source identification}\label{sec:source}

The scalar and bond operators must be checked separately. Raw I,
lines 23909--23918, transcribes CMP99b (3.18)--(3.19): its scalar
$Q'_n$ is a uniform block average with adjoint transports.
At flatness it is exactly $Q_0$. Raw I 23937--23944 transcribes
CMP99b (3.24)--(3.25) as \eqref{eq:source-R}. The orthogonal
projection calculation \eqref{eq:R-range} therefore identifies
the scalar gauge used above with the deposited homogeneous flat
source gauge. This is not the simpler projection $I-Q_0^*Q_0$.

Raw I 24349--24373 transcribes the corrected one-step average
of CMP98 (122)--(125). At the trivial configuration, all transports
are identities, $Y=Y_x=0$, and the four displayed correction
brackets vanish. The remaining linear operator is exactly
\[
 (QA)(c)=\sum_{x\in B(c_-)}L^{-d-1}A([x,x'])=Q_L^{\rm st}A(c).
\]
Raw I 47046--47048 gives the multilevel $Q_j$ as compositions of
these one-step operators. The centered digit and longitudinal-string
bijections of Notebook 06 give their flat composition
$Q_{L^j}^{\rm st}$. This is the corrected fine-to-coarse constraint,
distinct from the rooted next-step $Q^\flat$ used to eliminate
Gaussian pivot bonds.

For this homogeneous specialization, a source minimizing extension
with $QH=I$ and $RD_f^*H=0$ must equal $\HH_M$ by the uniqueness
argument above. If source notation instead reads $QH_0=\ell^{-1}$,
the normalized operator is $H_0\ell$, as recorded in Raw I
13638--13655. Omitting this normalization would compare different
maps. The conclusion identifies the explicit deposited flat
operator definitions and their stated minimizing property; it
does not reproduce the original nonlinear minimizer theorem.
The unavailable primary CMP99b text has not been independently
verified in this check.

An equivalent algebraic correction helps compare gauges. If $H$
is any minimizing right inverse with $QD_f=d_cQ_0$, then
\[
 H_{\rm src}=(I-D_fG'RD_f^*)H
\]
is feasible, has the same curl energy, and satisfies the source
gauge. Indeed $Q_0G'R=0$ and $\Delta G'R=R$ follow from
\eqref{eq:source-R}. Thus $QH_{\rm src}=I$ and
$RD_f^*H_{\rm src}=0$. Uniqueness identifies this correction
with the explicit formulas proved local above; locality is supplied
by those formulas, not assumed for a generic correction operator.

\section{What the localized extension can support}

The result supplies uniformly local fine potential, first-gradient
and curl legs for the homogeneous flat source operator. For example,
in a real differentiable family of self-adjoint stationary quadratic
forms satisfying
\[
 QH=I,\qquad \Delta H=Q^*q,\qquad q=H^*\Delta H,
\]
differentiation gives the exact identity
\[
 \dot q=H^*\dot\Delta H-H^*\dot Q^*q-q\dot QH.
\]
It uses $Q\dot H=-\dot QH$ and requires no separate derivative
bound on $\dot H$ in this first variation. Local variations of
these stated coefficients can consequently be tested against the
local flat legs just constructed.

This conditional identity by itself does not show that the deposited
independent-current family has this stationary representation, or
that its coefficient variations have the required weighted bounds.
No arbitrary nonflat background, hole, frozen boundary or mixed
history is covered here. No sharp cutoff is moved by an untracked
gauge change. The full complex carrier domain remains a separate
construction task. Notebook 09 now proves the stationary
representation and a uniform complex current neighborhood for the
homogeneous auxiliary family $G_\pi^J$, including all-order localized
current responses. Notebook 11 separately controls the additional
nonlinear constraint term and the full flat $G_1^J$ family in the
physical-cell algebra along $M=L^j$. Notebook 10 uses the full inverse to prove
locality of the projection onto feasible curvature variations and
an all-order expansion for small curvature-metric perturbations at
fixed averaging constraint.

\begin{thebibliography}{9}
\bibitem{repo}
\href{https://github.com/ethanwsussman/Yang--Mills-LLM-experiment/tree/e392ef45013a581664bab76b90867170b30fe10e}
{Pinned Yang--Mills-LLM-experiment account}, commit
\texttt{e392ef4501\allowbreak 3a581664ba\allowbreak b76b908671\allowbreak 70b30fe10e}.
\bibitem{balaban109} T. Ba\l aban,
\emph{Renormalization group approach to lattice gauge field theories I},
Commun. Math. Phys. 109 (1987), 249--301.
\href{https://scispace.com/pdf/renormalization-group-approach-to-lattice-gauge-field-3ygai0vxt1.pdf}
{Copy of the original article}.
\end{thebibliography}
\end{document}
